% Part: first-order-logic
% Chapter: completeness
% Section: complete-sets

% Definition of complete consistent sets. Properties of complete sets required
% for completeness proved are in provability.tex in the
% chapter on the proof system used.

\documentclass[../../../include/open-logic-section]{subfiles}

\begin{document}

\iftag{FOL}
      {\olfileid{fol}{com}{ccs}}
      {\olfileid{pl}{com}{ccs}}

\olsection{\usetoken{P}{sentence}ల సంపూర్ణ అవైరుధ్య సమితులు}

\begin{defn}[సంపూర్ణ సమితి]
\ollabel{def:complete-set} \tetoken{వాక్యాలు}{!!{sentence}s} సమితి~$\Gamma$ను
\emph{\tetoken{సంపూర్ణ}{!!{complete}}} అంటాం, అప్పుడు మాత్రమే: ఏ \tetoken{వాక్యం}{!!{sentence}}~$!A$కైనా
$!A \in \Gamma$ లేదా $\lnot !A \in \Gamma$.
\end{defn}

\begin{explain}
\tetoken{సంపూర్ణ}{!!^{complete}} వాక్యాల సమితులు ఏ ప్రశ్ననూ సమాధానం లేకుండా వదలవు. ఏ
\tetoken{వాక్యం}{!!{sentence}}~$!A$కైనా, $!A$ సత్యమా అసత్యమా అనేది~$\Gamma$
``చెబుతుంది.'' \tetoken{సంపూర్ణ}{!!{complete}} సమితుల ప్రాముఖ్యం సంపూర్ణతా సిద్ధాంతపు
నిరూపణకే పరిమితం కాదు. ఉదాహరణకు, \tetoken{సంపూర్ణంగా}{!!{complete}}, స్వీకృతీకరించదగినదిగా
ఉన్న సిద్ధాంతం ఎల్లప్పుడూ నిర్ణయించదగినదే.
\end{explain}

\begin{explain}
సంపూర్ణతా నిరూపణలో \tetoken{సంపూర్ణ}{!!^{complete}} అవైరుధ్య సమితులు ముఖ్యమైనవి; ఎందుకంటే
ప్రతి అవైరుధ్య \tetoken{వాక్యాలు}{!!{sentence}s} సమితి~$\Gamma$ ఏదో ఒక \tetoken{సంపూర్ణ}{!!a{complete}}
అవైరుధ్య సమితి~$\Gamma^*$లో ఇమిడి ఉంటుందని హామీ ఇవ్వగలం.
\tetoken{ఒక సంపూర్ణ}{!!^a{complete}} అవైరుధ్య సమితిలో ప్రతి \tetoken{వాక్యం}{!!{sentence}}~$!A$కు $!A$ లేదా
దాని నిషేధం $\lnot !A$ ఉంటుంది, కానీ రెండూ ఉండవు. ప్రత్యేకంగా ఇది
\iftag{FOL}{పరమాణు \tetoken{వాక్యాలు}{!!{sentence}s}}{ప్రతిజ్ఞావాక్య చరాలు}కు వర్తిస్తుంది.
కాబట్టి \iftag{FOL}{\tetoken{స్థిరసంకేతాలతో}{!!{constant}s} తగినట్లు విస్తరించిన భాషలోని}{}
\tetoken{ఒక సంపూర్ణ}{!!a{complete}} అవైరుధ్య సమితి నుంచి, ఏ
\iftag{FOL}{పరమాణు \tetoken{వాక్యాలు}{!!{sentence}s}}{\tetoken{ప్రతిజ్ఞావాక్య చరాలు}{!!{propositional variable}s}}
~$\Gamma^*$లో ఉన్నాయో దాని ప్రకారం
\iftag{FOL}{\tetoken{విధేయ సంకేతాలు}{!!{predicate}s} అర్థవివరణ}{ప్రతిజ్ఞావాక్య చరాలకు నిర్దేశించే
సత్యమూల్యం} నిర్వచించబడే \iftag{FOL}{\tetoken{ఒక నిర్మాణం}{!!a{structure}}}{\tetoken{ఒక సత్యమూల్య కేటాయింపు}{!!a{valuation}}}ను
నిర్మించవచ్చు. ఈ \iftag{FOL}{\tetoken{నిర్మాణం}{!!{structure}}}{\tetoken{సత్యమూల్య కేటాయింపు}{!!{valuation}}} తరువాత
~$\Gamma^*$లోని అన్ని \tetoken{వాక్యాలను}{!!{sentence}s} (అందువల్ల~$\Gamma$లోని
వాటన్నింటినీ కూడా) సత్యం చేస్తుందని చూపవచ్చు. ఈ చివరి వాస్తవపు
నిరూపణకు $\lnot !A \in \Gamma^*$ అయినప్పుడు, అప్పుడు మాత్రమే $!A
\notin \Gamma^*$; $(!A \lor !B) \in \Gamma^*$ అయినప్పుడు, అప్పుడు
మాత్రమే $!A \in \Gamma^*$ లేదా $!B \in \Gamma^*$; తదితర ధర్మాలు
కావాలి.
\end{explain}

కింద తరచుగా $\Proves$ యొక్క స్వావర్తనత్వం, ఏకదిశత, సంక్రామకత్వ
ధర్మాలను అవ్యక్తంగా ఉపయోగిస్తాం (చూడండి
\iftag{FOL}{%
  \tagrefs{prfSC/{fol:seq:ptn:sec},prfND/{fol:ntd:ptn:sec},prfAX/{fol:axd:ptn:sec},prfTab/{fol:tab:ptn:sec}}}{%
  \tagrefs{prfSC/{pl:seq:ptn:sec},prfND/{pl:ntd:ptn:sec},prfAX/{pl:axd:ptn:sec},prfTab/{pl:tab:ptn:sec}}}).

\begin{prop}
\ollabel{prop:ccs}
$\Gamma$ \tetoken{సంపూర్ణ}{!!{complete}}, అవైరుధ్యమైనది అనుకుందాం. అప్పుడు:
\begin{enumerate}
\item \ollabel{prop:ccs-prov-in} $\Gamma \Proves !A$ అయితే, $!A \in
  \Gamma$.

\tagitem{prvAnd}{\ollabel{prop:ccs-and} $!A \land !B \in \Gamma$
  అయినప్పుడు, అప్పుడు మాత్రమే $!A \in \Gamma$, $!B \in \Gamma$
  రెండూ.}{}

\tagitem{prvOr}{\ollabel{prop:ccs-or} $!A \lor !B \in \Gamma$
  అయినప్పుడు, అప్పుడు మాత్రమే $!A \in \Gamma$ లేదా $!B \in \Gamma$.}{}

\tagitem{prvIf}{\ollabel{prop:ccs-if} $!A \lif !B \in \Gamma$
  అయినప్పుడు, అప్పుడు మాత్రమే $!A \notin \Gamma$ లేదా $!B \in
  \Gamma$.}{}
\end{enumerate}
\end{prop}

\begin{proof}
కింద అన్నింటిలోనూ~$\Gamma$ \tetoken{సంపూర్ణ}{!!{complete}}, అవైరుధ్యమైనది అని
అనుకుందాం.
\begin{enumerate}
\item $\Gamma \Proves !A$ అయితే, $!A \in \Gamma$.

$\Gamma \Proves !A$ అని అనుకుందాం. దీనికి విరుద్ధంగా $!A \notin
\Gamma$ అని అనుకుందాం. $\Gamma$ \tetoken{సంపూర్ణ}{!!{complete}} కనుక $\lnot !A \in
\Gamma$. కింది ఫలితం ప్రకారం
\iftag{FOL}{%
  \tagrefs{prfSC/{fol:seq:prv:prop:explicit-inc},
    prfND/{fol:ntd:prv:prop:explicit-inc},
    prfAX/{fol:axd:prv:prop:explicit-inc},
    prfTab/{fol:tab:prv:prop:explicit-inc}}}{%
  \tagrefs{prfSC/{pl:seq:prv:prop:explicit-inc},
    prfND/{pl:ntd:prv:prop:explicit-inc},
    prfAX/{pl:axd:prv:prop:explicit-inc},
    prfTab/{pl:tab:prv:prop:explicit-inc}}},
$\Gamma$ వైరుధ్యభరితమవుతుంది. ఇది~$\Gamma$ అవైరుధ్యమనే పరికల్పనకు
విరుద్ధం. కాబట్టి $!A \notin \Gamma$ కావడం అసాధ్యం; అందువల్ల $!A
\in \Gamma$.

\tagitem{defAnd}{}{%
\iftag{probAnd}{అభ్యాసం.}{%
$!A \land !B \in \Gamma$ అయినప్పుడు, అప్పుడు మాత్రమే $!A \in
\Gamma$, $!B \in \Gamma$ రెండూ:

ముందుకు వెళ్లే దిశలో $!A \land !B \in \Gamma$ అని అనుకుందాం. అప్పుడు
కింది ఫలితంలోని అంశం~(1) ప్రకారం
\iftag{FOL}{%
  \tagrefs{prfSC/{fol:seq:ppr:prop:provability-land},
    prfND/{fol:ntd:ppr:prop:provability-land},
    prfAX/{fol:axd:ppr:prop:provability-land},
    prfTab/{fol:tab:ppr:prop:provability-land}}}{%
  \tagrefs{prfSC/{pl:seq:ppr:prop:provability-land},
    prfND/{pl:ntd:ppr:prop:provability-land},
    prfAX/{pl:axd:ppr:prop:provability-land},
    prfTab/{pl:tab:ppr:prop:provability-land}}},
$\Gamma \Proves !A$, $\Gamma \Proves !B$. \olref{prop:ccs-prov-in}
ప్రకారం $!A \in \Gamma$, $!B \in \Gamma$; కావలసింది ఇదే.

తిరుగు దిశలో $!A \in \Gamma$, $!B \in \Gamma$ అనుకుందాం. కింది
ఫలితంలోని అంశం~(2) ప్రకారం
\iftag{FOL}{%
  \tagrefs{prfSC/{fol:seq:ppr:prop:provability-land},
    prfND/{fol:ntd:ppr:prop:provability-land},
    prfAX/{fol:axd:ppr:prop:provability-land},
    prfTab/{fol:tab:ppr:prop:provability-land}}}{%
  \tagrefs{prfSC/{pl:seq:ppr:prop:provability-land},
    prfND/{pl:ntd:ppr:prop:provability-land},
    prfAX/{pl:axd:ppr:prop:provability-land},
    prfTab/{pl:tab:ppr:prop:provability-land}}},
$\Gamma \Proves !A \land !B$. \olref{prop:ccs-prov-in} ప్రకారం $!A
\land !B \in \Gamma$.}}

\tagitem{defOr}{}{%
\iftag{probOr}{అభ్యాసం.}{%
మొదట, $!A \lor !B \in \Gamma$ అయితే $!A \in \Gamma$ లేదా $!B \in
\Gamma$ అని చూపుతాం. $!A \lor !B \in \Gamma$ అయినా $!A \notin
\Gamma$, $!B \notin \Gamma$ అని అనుకుందాం. $\Gamma$ \tetoken{సంపూర్ణ}{!!{complete}}
కనుక $\lnot !A \in \Gamma$, $\lnot !B \in \Gamma$. కింది ఫలితంలోని
అంశం~(1) ప్రకారం
\iftag{FOL}{%
  \tagrefs{prfSC/{fol:seq:ppr:prop:provability-lor},
    prfND/{fol:ntd:ppr:prop:provability-lor},
    prfAX/{fol:axd:ppr:prop:provability-lor},
    prfTab/{fol:tab:ppr:prop:provability-lor}}}{%
  \tagrefs{prfSC/{pl:seq:ppr:prop:provability-lor},
    prfND/{pl:ntd:ppr:prop:provability-lor},
    prfAX/{pl:axd:ppr:prop:provability-lor},
    prfTab/{pl:tab:ppr:prop:provability-lor}}},
$\Gamma$ వైరుధ్యభరితమవుతుంది; ఇది వైరుధ్యం. కాబట్టి $!A \in
\Gamma$ లేదా $!B \in \Gamma$.

తిరుగు దిశలో $!A \in \Gamma$ లేదా $!B \in \Gamma$ అనుకుందాం. కింది
ఫలితంలోని అంశం~(2) ప్రకారం
\iftag{FOL}{%
  \tagrefs{prfSC/{fol:seq:ppr:prop:provability-lor},
    prfND/{fol:ntd:ppr:prop:provability-lor},
    prfAX/{fol:axd:ppr:prop:provability-lor},
    prfTab/{fol:tab:ppr:prop:provability-lor}}}{%
  \tagrefs{prfSC/{pl:seq:ppr:prop:provability-lor},
    prfND/{pl:ntd:ppr:prop:provability-lor},
    prfAX/{pl:axd:ppr:prop:provability-lor},
    prfTab/{pl:tab:ppr:prop:provability-lor}}},
$\Gamma \Proves !A \lor !B$. \olref{prop:ccs-prov-in} ప్రకారం $!A
\lor !B \in \Gamma$; కావలసింది ఇదే.}}

\tagitem{defIf}{}{%
\iftag{probIf}{అభ్యాసం.}{%
ముందుకు వెళ్లే దిశలో $!A \lif !B \in \Gamma$ అని అనుకుందాం; దీనికి
విరుద్ధంగా $!A \in \Gamma$, $!B \notin \Gamma$ అని అనుకుందాం. ఈ
పరికల్పనలపై $!A \lif !B \in \Gamma$, $!A \in \Gamma$. కింది
ఫలితంలోని అంశం~(1) ప్రకారం
\iftag{FOL}{%
  \tagrefs{prfSC/{fol:seq:ppr:prop:provability-lif},
    prfND/{fol:ntd:ppr:prop:provability-lif},
    prfAX/{fol:axd:ppr:prop:provability-lif},
    prfTab/{fol:tab:ppr:prop:provability-lif}}}{%
  \tagrefs{prfSC/{pl:seq:ppr:prop:provability-lif},
    prfND/{pl:ntd:ppr:prop:provability-lif},
    prfAX/{pl:axd:ppr:prop:provability-lif},
    prfTab/{pl:tab:ppr:prop:provability-lif}}},
$\Gamma \Proves !B$. కానీ \olref{prop:ccs-prov-in} ప్రకారం అప్పుడు
$!B \in \Gamma$; ఇది $!B \notin \Gamma$ అనే పరికల్పనకు విరుద్ధం.

తిరుగు దిశలో మొదట $!A \notin \Gamma$ అయ్యే సందర్భాన్ని పరిశీలిద్దాం.
$\Gamma$ \tetoken{సంపూర్ణ}{!!{complete}} కనుక $\lnot !A \in \Gamma$. కింది ఫలితంలోని
అంశం~(2) ప్రకారం
\iftag{FOL}{%
  \tagrefs{prfSC/{fol:seq:ppr:prop:provability-lif},
    prfND/{fol:ntd:ppr:prop:provability-lif},
    prfAX/{fol:axd:ppr:prop:provability-lif},
    prfTab/{fol:tab:ppr:prop:provability-lif}}}{%
  \tagrefs{prfSC/{pl:seq:ppr:prop:provability-lif},
    prfND/{pl:ntd:ppr:prop:provability-lif},
    prfAX/{pl:axd:ppr:prop:provability-lif},
    prfTab/{pl:tab:ppr:prop:provability-lif}}},
$\Gamma \Proves !A \lif !B$. మళ్లీ \olref{prop:ccs-prov-in} ద్వారా
$!A \lif !B \in \Gamma$ వస్తుంది; కావలసింది ఇదే.

ఇప్పుడు $!B \in \Gamma$ అయ్యే సందర్భాన్ని పరిశీలిద్దాం. కింది
ఫలితంలోని అంశం~(2) ప్రకారం
\iftag{FOL}{%
  \tagrefs{prfSC/{fol:seq:ppr:prop:provability-lif},
    prfND/{fol:ntd:ppr:prop:provability-lif},
    prfTab/{fol:tab:ppr:prop:provability-lif},
    prfAX/{fol:axd:ppr:prop:provability-lif}}}{%
  \tagrefs{prfSC/{pl:seq:ppr:prop:provability-lif},
    prfND/{pl:ntd:ppr:prop:provability-lif},
    prfTab/{pl:tab:ppr:prop:provability-lif},
    prfAX/{pl:axd:ppr:prop:provability-lif}}},
మళ్లీ $\Gamma \Proves !A \lif !B$. \olref{prop:ccs-prov-in} ప్రకారం
$!A \lif !B \in \Gamma$.}}
\end{enumerate}
\end{proof}

\tagprob{FOL}
\begin{prob}
\olref[fol][com][ccs]{prop:ccs} నిరూపణను పూర్తి చేయండి.
\end{prob}
\tagendprob

\tagprob{notFOL}
\begin{prob}
\olref[pl][com][ccs]{prop:ccs} నిరూపణను పూర్తి చేయండి.
\end{prob}
\tagendprob

\end{document}
